% ============================================================
\chapter{Pr\'ediction Bay\'esienne}
\label{ch:prediction}
% ============================================================

En statistique fr\'equentiste, la pr\'ediction utilise un estimateur ponctuel~: on estime $\theta$, puis on pr\'edit $\tilde{y}$ comme si $\hat\theta$ \'etait la vraie valeur. Cette approche ignore l'incertitude sur le param\`etre et produit des intervalles de pr\'ediction trop \'etroits. L'approche bay\'esienne est fondamentalement diff\'erente~: au lieu de fixer $\theta$, on \emph{int\`egre} sur toutes les valeurs possibles de $\theta$, pond\'er\'ees par la loi a posteriori. La distribution pr\'edictive qui en r\'esulte incorpore naturellement deux sources d'incertitude --- l'al\'ea intrins\`eque des donn\'ees et l'incertitude sur le mod\`ele --- et produit des pr\'evisions mieux calibr\'ees.

\begin{intuition}[title=\textbf{Id\'ee centrale}]
La pr\'ediction bay\'esienne int\`egre l'incertitude sur le param\`etre
en moyennant la distribution pr\'edictive sur le posterior. Cela produit
des intervalles pr\'edictifs calibr\'es et \'evite le sur-ajustement
inh\'erent \`a l'utilisation d'estimations ponctuelles.
\end{intuition}

% ------------------------------------------------------------
\section{Distribution pr\'edictive a priori}
% ------------------------------------------------------------

\begin{definition}[Distribution pr\'edictive a priori]
La \textbf{distribution pr\'edictive a priori} (ou marginale) d'une
nouvelle observation $\tilde{x}$ est~:
\[
  m(\tilde{x}) = \int_\Theta f(\tilde{x} \mid \theta)\,\pi(\theta)\,d\theta.
\]
Elle ne d\'epend pas des donn\'ees et repr\'esente la pr\'ediction avant
toute observation.
\end{definition}

\begin{exemple}
\textbf{Mod\`ele Bernoulli--B\^eta.}
Avec $\theta \sim \text{Be}(a, b)$ et $\tilde{X} \mid \theta \sim \text{Ber}(\theta)$~:
\[
  m(\tilde{x} = 1) = \E[\theta] = \frac{a}{a + b}.
\]
Pour le prior de Laplace $\text{Be}(1, 1)$~: $m(\tilde{x} = 1) = 1/2$.
\end{exemple}

\begin{exemple}
\textbf{Mod\`ele Normal--Normal.}
Avec $\mu \sim \mathcal{N}(\mu_0, \tau_0^2)$ et
$\tilde{X} \mid \mu \sim \mathcal{N}(\mu, \sigma^2)$~:
\[
  m(\tilde{x}) = \int \mathcal{N}(\tilde{x} \mid \mu, \sigma^2)
  \cdot \mathcal{N}(\mu \mid \mu_0, \tau_0^2)\,d\mu
  = \mathcal{N}(\tilde{x} \mid \mu_0, \sigma^2 + \tau_0^2).
\]
La variance pr\'edictive $\sigma^2 + \tau_0^2$ additionne l'incertitude
al\'eatoire et l'incertitude \'epist\'emique.
\end{exemple}

% ------------------------------------------------------------
\section{Distribution pr\'edictive a posteriori}
% ------------------------------------------------------------

\begin{definition}[Distribution pr\'edictive a posteriori]
La \textbf{distribution pr\'edictive a posteriori} de $\tilde{x}$
\'etant donn\'e $x = (x_1, \ldots, x_n)$ est~:
\[
  f(\tilde{x} \mid x) = \int_\Theta f(\tilde{x} \mid \theta)\,\pi(\theta \mid x)\,d\theta.
\]
\end{definition}

\begin{formules}[title=\textbf{D\'ecomposition de la variance pr\'edictive}]
\[
  \Var[\tilde{X} \mid x]
  = \underbrace{\E_{\theta \mid x}[\Var[\tilde{X} \mid \theta]]}_{\text{incertitude al\'eatoire}}
  + \underbrace{\Var_{\theta \mid x}[\E[\tilde{X} \mid \theta]]}_{\text{incertitude \'epist\'emique}}.
\]
Cette d\'ecomposition est la \textbf{loi de la variance totale}.
\end{formules}

\begin{theoreme}[Pr\'edictive a posteriori Normal--Normal]
Si $X_i \overset{\text{iid}}{\sim} \mathcal{N}(\mu, \sigma^2)$ avec
$\mu \mid x \sim \mathcal{N}(\mu_n, \tau_n^2)$, alors~:
\[
  \tilde{X} \mid x \sim \mathcal{N}(\mu_n, \sigma^2 + \tau_n^2).
\]
\end{theoreme}

\begin{proof}
Conditionnellement \`a $\mu$, $\tilde{X} \sim \mathcal{N}(\mu, \sigma^2)$.
En int\'egrant sur $\mu \mid x \sim \mathcal{N}(\mu_n, \tau_n^2)$,
la somme de deux gaussiennes ind\'ependantes donne~:
\[
  \tilde{X} \mid x \sim \mathcal{N}(\mu_n, \sigma^2 + \tau_n^2).
\]
\end{proof}

\begin{theoreme}[Pr\'edictive a posteriori Bernoulli--B\^eta]
Si $X_i \overset{\text{iid}}{\sim} \text{Ber}(\theta)$ avec
$\theta \mid x \sim \text{Be}(a_n, b_n)$, alors~:
\[
  \PP(\tilde{X} = 1 \mid x) = \E[\theta \mid x] = \frac{a_n}{a_n + b_n}.
\]
C'est la \textbf{r\`egle de succession de Laplace} quand $a = b = 1$~:
$\PP(\tilde{X} = 1 \mid x) = (s+1)/(n+2)$.
\end{theoreme}

\begin{theoreme}[Pr\'edictive Poisson--Gamma]
Si $X_i \overset{\text{iid}}{\sim} \text{Poi}(\lambda)$ avec
$\lambda \mid x \sim \text{Ga}(\alpha_n, \beta_n)$, la pr\'edictive est
une \textbf{binomiale n\'egative}~:
\[
  \tilde{X} \mid x \sim \text{NegBin}\!\left(\alpha_n, \frac{\beta_n}{\beta_n + 1}\right).
\]
\end{theoreme}

\begin{proof}
\begin{align*}
  \PP(\tilde{X} = k \mid x)
  &= \int_0^\infty \frac{e^{-\lambda}\lambda^k}{k!}
  \cdot \frac{\beta_n^{\alpha_n}}{\Gamma(\alpha_n)}
  \lambda^{\alpha_n - 1} e^{-\beta_n \lambda}\,d\lambda \\
  &= \frac{\beta_n^{\alpha_n}}{k!\,\Gamma(\alpha_n)}
  \int_0^\infty \lambda^{k + \alpha_n - 1} e^{-(\beta_n + 1)\lambda}\,d\lambda \\
  &= \frac{\beta_n^{\alpha_n}}{k!\,\Gamma(\alpha_n)}
  \cdot \frac{\Gamma(k + \alpha_n)}{(\beta_n + 1)^{k + \alpha_n}} \\
  &= \binom{k + \alpha_n - 1}{k}
  \left(\frac{\beta_n}{\beta_n + 1}\right)^{\alpha_n}
  \left(\frac{1}{\beta_n + 1}\right)^k.
\end{align*}
\end{proof}

\begin{theoreme}[Pr\'edictive Normal--NIG]
Si $(\mu, \sigma^2) \mid x \sim \text{NIG}(\mu_n, \kappa_n, \alpha_n, \beta_n)$,
la pr\'edictive a posteriori est une Student~:
\[
  \tilde{X} \mid x \sim t_{2\alpha_n}\!\left(\mu_n,\;
  \frac{\beta_n(\kappa_n + 1)}{\alpha_n \kappa_n}\right).
\]
\end{theoreme}

% ------------------------------------------------------------
\section{V\'erification pr\'edictive a posteriori}
% ------------------------------------------------------------

\begin{definition}[Posterior predictive check]
La \textbf{v\'erification pr\'edictive a posteriori} consiste \`a
simuler des r\'epliques $x^{\text{rep}}$ depuis la pr\'edictive et \`a
les comparer aux donn\'ees observ\'ees pour \'evaluer l'ad\'equation du
mod\`ele.
\end{definition}

\begin{definition}[$p$-valeur pr\'edictive bay\'esienne]
La \textbf{$p$-valeur pr\'edictive} pour une statistique de test
$T(x)$ est~:
\[
  p_B = \PP(T(x^{\text{rep}}) \geq T(x_{\text{obs}}) \mid x_{\text{obs}}).
\]
Des valeurs proches de $0$ ou $1$ sugg\`erent un mauvais ajustement.
\end{definition}

\begin{algorithme}[title=\textbf{V\'erification pr\'edictive a posteriori}]
\begin{enumerate}
  \item Pour $t = 1, \ldots, T$~:
  \begin{enumerate}
    \item Tirer $\theta^{(t)} \sim \pi(\theta \mid x_{\text{obs}})$ (depuis MCMC).
    \item Tirer $x^{\text{rep},(t)} \sim f(x \mid \theta^{(t)})$.
    \item Calculer $T^{(t)} = T(x^{\text{rep},(t)})$.
  \end{enumerate}
  \item Estimer $p_B = \frac{1}{T}\sum_{t=1}^T \mathbf{1}(T^{(t)} \geq T(x_{\text{obs}}))$.
  \item Tracer la distribution de $T^{(t)}$ et comparer \`a $T(x_{\text{obs}})$.
\end{enumerate}
\end{algorithme}

% ------------------------------------------------------------
\section{Intervalles pr\'edictifs}
% ------------------------------------------------------------

\begin{definition}[Intervalle pr\'edictif]
Un \textbf{intervalle pr\'edictif} \`a $100(1-\alpha)\%$ pour $\tilde{X}$ est
un intervalle $[a, b]$ tel que~:
\[
  \PP(\tilde{X} \in [a, b] \mid x) = 1 - \alpha.
\]
\end{definition}

\begin{remarque}
L'intervalle pr\'edictif est toujours plus large que l'intervalle de
cr\'edibilit\'e pour $\theta$, car il int\`egre \`a la fois l'incertitude
param\'etrique et l'al\'ea de la nouvelle observation.
\end{remarque}

% ------------------------------------------------------------
\section{Pr\'ediction bay\'esienne vs. plug-in}
% ------------------------------------------------------------

\begin{attention}[title=\textbf{Danger du plug-in}]
L'approche \textit{plug-in} remplace $\theta$ par une estimation
ponctuelle $\hat{\theta}$ et pr\'edit via $f(\tilde{x} \mid \hat{\theta})$.
Elle sous-estime l'incertitude pr\'edictive car elle ignore l'incertitude
sur $\theta$. La pr\'ediction bay\'esienne est mieux calibr\'ee.
\end{attention}

% ------------------------------------------------------------
\section{Impl\'ementation Python}
% ------------------------------------------------------------

\begin{codeblock}[title=\textbf{Distribution pr\'edictive a posteriori}]
\begin{minted}{python}
import numpy as np
import pymc as pm
import arviz as az
import matplotlib.pyplot as plt

# Données gaussiennes
np.random.seed(42)
n = 30
mu_true, sigma_true = 5.0, 2.0
data = np.random.normal(mu_true, sigma_true, n)

# Modèle bayésien
with pm.Model() as model_pred:
    mu = pm.Normal("mu", mu=0, sigma=10)
    sigma = pm.HalfNormal("sigma", sigma=10)
    y = pm.Normal("y", mu=mu, sigma=sigma, observed=data)
    # Prédictive
    y_pred = pm.Normal("y_pred", mu=mu, sigma=sigma)
    trace = pm.sample(3000, random_seed=42)
    ppc = pm.sample_posterior_predictive(trace,
                                         random_seed=42)

# Visualisation
fig, axes = plt.subplots(1, 2, figsize=(12, 4))

# Prédictive vs données
az.plot_ppc(ppc, observed_adata=trace, ax=axes[0])
axes[0].set_title("Vérification prédictive")

# Intervalle prédictif
y_pred_samples = ppc.posterior_predictive["y_pred"].values.flatten()
ci_pred = np.percentile(y_pred_samples, [2.5, 97.5])
axes[1].hist(y_pred_samples, bins=50, density=True,
             alpha=0.5, label="Prédictive")
axes[1].axvline(ci_pred[0], color="red", linestyle="--",
                label=f"IC 95%: [{ci_pred[0]:.1f}, {ci_pred[1]:.1f}]")
axes[1].axvline(ci_pred[1], color="red", linestyle="--")
axes[1].legend()
axes[1].set_title("Distribution prédictive a posteriori")

plt.tight_layout()
plt.savefig("predictive.pdf")
\end{minted}
\end{codeblock}

\begin{codeblock}[title=\textbf{$p$-valeur pr\'edictive bay\'esienne}]
\begin{minted}{python}
# Statistique de test: écart-type des données
T_obs = data.std()

# Simuler des répliques
T_rep = []
mu_samples = trace.posterior["mu"].values.flatten()
sigma_samples = trace.posterior["sigma"].values.flatten()
for i in range(len(mu_samples)):
    x_rep = np.random.normal(mu_samples[i],
                              sigma_samples[i], n)
    T_rep.append(x_rep.std())
T_rep = np.array(T_rep)

# p-valeur prédictive
p_B = np.mean(T_rep >= T_obs)
print(f"T(x_obs) = {T_obs:.3f}")
print(f"p-valeur prédictive = {p_B:.3f}")

# Visualisation
plt.figure(figsize=(8, 4))
plt.hist(T_rep, bins=50, density=True, alpha=0.5,
         label="T(x_rep)")
plt.axvline(T_obs, color="red",
            label=f"T(x_obs) = {T_obs:.2f}")
plt.xlabel("Écart-type")
plt.ylabel("Densité")
plt.legend()
plt.title(f"p-valeur prédictive = {p_B:.3f}")
plt.tight_layout()
plt.savefig("ppc_pvalue.pdf")
\end{minted}
\end{codeblock}

\begin{codeblock}[title=\textbf{Pr\'edictive Poisson--Gamma (binomiale n\'egative)}]
\begin{minted}{python}
from scipy import stats

# Données Poisson
np.random.seed(123)
data_pois = np.random.poisson(3.0, size=20)

# Prior Gamma(2, 1)
alpha_n = 2 + data_pois.sum()
beta_n = 1 + len(data_pois)
p_nb = beta_n / (beta_n + 1)

# Prédictive: NegBin(alpha_n, p_nb)
x_grid = np.arange(0, 15)
pred_pmf = stats.nbinom.pmf(x_grid, alpha_n, p_nb)

plt.figure(figsize=(8, 4))
plt.bar(x_grid, pred_pmf, alpha=0.7,
        label="Prédictive NegBin")
plt.xlabel(r"$\tilde{x}$")
plt.ylabel("Probabilité")
plt.title("Prédictive Poisson-Gamma")
plt.legend()
plt.tight_layout()
plt.savefig("pred_negbin.pdf")
\end{minted}
\end{codeblock}

% ------------------------------------------------------------
\section{Exercices}
% ------------------------------------------------------------

\begin{exercice}
D\'eriver la distribution pr\'edictive a posteriori pour le mod\`ele
Bernoulli--B\^eta pour $m$ nouvelles observations (pas une seule).
Montrer qu'il s'agit d'une loi b\^eta-binomiale.
\end{exercice}

\begin{exercice}
Montrer que la variance pr\'edictive est toujours sup\'erieure ou
\'egale \`a la variance conditionnelle $\Var[\tilde{X} \mid \theta]$
pour toute valeur de $\theta$.
\end{exercice}

\begin{exercice}
D\'eriver la distribution pr\'edictive pour le mod\`ele Normal--NIG et
montrer qu'il s'agit d'une Student.
\end{exercice}

\begin{exercice}
\textbf{(Num\'erique)} Simuler des donn\'ees Poisson et comparer~:
\begin{enumerate}
  \item la pr\'edictive bay\'esienne (binomiale n\'egative),
  \item la pr\'edictive plug-in (Poisson avec $\hat{\lambda} = \bar{x}$).
\end{enumerate}
Calculer la couverture r\'eelle des intervalles pr\'edictifs \`a $95\%$
des deux approches sur $1000$ r\'ep\'etitions.
\end{exercice}

\begin{exercice}
Effectuer une v\'erification pr\'edictive a posteriori pour un mod\`ele
Poisson appliqu\'e \`a des donn\'ees sur-dispers\'ees. Utiliser la
variance comme statistique de test.
\end{exercice}
